% Options for packages loaded elsewhere
\PassOptionsToPackage{unicode}{hyperref}
\PassOptionsToPackage{hyphens}{url}
\documentclass[
]{article}
\usepackage{xcolor}
\usepackage{amsmath,amssymb}
\setcounter{secnumdepth}{5}
\usepackage{iftex}
\ifPDFTeX
  \usepackage[T1]{fontenc}
  \usepackage[utf8]{inputenc}
  \usepackage{textcomp} % provide euro and other symbols
\else % if luatex or xetex
  \usepackage{unicode-math} % this also loads fontspec
  \defaultfontfeatures{Scale=MatchLowercase}
  \defaultfontfeatures[\rmfamily]{Ligatures=TeX,Scale=1}
\fi
% \usepackage{lmodern}
\ifPDFTeX\else
  % xetex/luatex font selection
\fi
% Use upquote if available, for straight quotes in verbatim environments
\IfFileExists{upquote.sty}{\usepackage{upquote}}{}
\IfFileExists{microtype.sty}{% use microtype if available
  \usepackage[]{microtype}
  \UseMicrotypeSet[protrusion]{basicmath} % disable protrusion for tt fonts
}{}
\makeatletter
\@ifundefined{KOMAClassName}{% if non-KOMA class
  \IfFileExists{parskip.sty}{%
    \usepackage{parskip}
  }{% else
    \setlength{\parindent}{0pt}
    \setlength{\parskip}{6pt plus 2pt minus 1pt}}
}{% if KOMA class
  \KOMAoptions{parskip=half}}
\makeatother
\usepackage{longtable,booktabs,array}
\newcounter{none} % for unnumbered tables
\usepackage{calc} % for calculating minipage widths
% Correct order of tables after \paragraph or \subparagraph
\usepackage{etoolbox}
\makeatletter
\patchcmd\longtable{\par}{\if@noskipsec\mbox{}\fi\par}{}{}
\makeatother
% Allow footnotes in longtable head/foot
\IfFileExists{footnotehyper.sty}{\usepackage{footnotehyper}}{\usepackage{footnote}}
\makesavenoteenv{longtable}
\setlength{\emergencystretch}{3em} % prevent overfull lines
\providecommand{\tightlist}{%
  \setlength{\itemsep}{0pt}\setlength{\parskip}{0pt}}
\usepackage[]{natbib}
\bibliographystyle{plainnat}
\makeatletter
\@ifpackageloaded{subcaption}{}{\usepackage{subcaption}}
\@ifpackageloaded{caption}{}{\usepackage{caption}}
\captionsetup[subfigure]{margin=0.5em}
\AtBeginDocument{%
\renewcommand*\figurename{Figure}
\renewcommand*\tablename{Table}
}
\AtBeginDocument{%
\renewcommand*\listfigurename{List of Figures}
\renewcommand*\listtablename{List of Tables}
}
\newcounter{pandoccrossref@subfigures@footnote@counter}
\newenvironment{pandoccrossrefsubfigures}{%
\setcounter{pandoccrossref@subfigures@footnote@counter}{0}
\begin{figure}\centering%
\gdef\global@pandoccrossref@subfigures@footnotes{}%
\DeclareRobustCommand{\footnote}[1]{\footnotemark%
\stepcounter{pandoccrossref@subfigures@footnote@counter}%
\ifx\global@pandoccrossref@subfigures@footnotes\empty%
\gdef\global@pandoccrossref@subfigures@footnotes{{##1}}%
\else%
\g@addto@macro\global@pandoccrossref@subfigures@footnotes{, {##1}}%
\fi}}%
{\end{figure}%
\addtocounter{footnote}{-\value{pandoccrossref@subfigures@footnote@counter}}
\@for\f:=\global@pandoccrossref@subfigures@footnotes\do{\stepcounter{footnote}\footnotetext{\f}}%
\gdef\global@pandoccrossref@subfigures@footnotes{}}
\@ifpackageloaded{float}{}{\usepackage{float}}
\floatstyle{ruled}
\@ifundefined{c@chapter}{\newfloat{codelisting}{h}{lop}}{\newfloat{codelisting}{h}{lop}[chapter]}
\floatname{codelisting}{Listing}
\newcommand*\listoflistings{\listof{codelisting}{List of Listings}}
\makeatother
\usepackage{bookmark}
\IfFileExists{xurl.sty}{\usepackage{xurl}}{} % add URL line breaks if available
\urlstyle{same}
% fallback for those not using the hyperref driver hyperxmp:
\makeatletter
\@ifundefined{xmpquote}{\newcommand{\xmpquote}[1]{#1}}{}
\makeatother
\hypersetup{
  colorlinks=true,linkcolor=red,urlcolor=red,citecolor=red,anchorcolor=red,filecolor=red,
  pdfcreator={LaTeX via pandoc}}

\author{}
\date{}

\title{Redacted Report Template: Disclosure Control and Release Audit}
\author{Research Template Author\\\footnotesize{\texttt{author@research-template.org}}}
\date{2026-07-09}

\usepackage[left=0.16in,right=0.16in,top=0.48in,bottom=0.52in]{geometry}
\usepackage{xcolor}
\usepackage{graphicx}
\definecolor{redactedReportPage}{HTML}{FFFFFF}
\pagecolor{redactedReportPage}
\color{black}
\setlength{\parindent}{0pt}
\setlength{\parskip}{0.45em}

\begin{document}

\begin{titlepage}
\centering
\vspace*{2cm}
{\LARGE\bfseries Redacted Report Template: Disclosure Control and Release Audit\par}
\vspace{0.75em}
{\large A public exemplar for sanitized release workflows\par}
\vfill
\makeatletter
{\@author\par}
\vspace{1em}
{\@date\par}
\makeatother
\vfill
\end{titlepage}
\thispagestyle{empty}
\tableofcontents
\newpage


\section{Abstract}\label{abstract}

This exemplar demonstrates redaction ledgers, classification ceilings,
source-protection checks, release authority, and mosaic-risk review for
sanitized public reports.

\newpage

\section{Methods}\label{methods}

The method audits each segment against a public classification ceiling,
validates redaction spans, checks orphan decisions, checks
source-control coverage, normalizes classification taxonomies, exports
sanitized release packets, builds source-safe redaction ledgers, records
segment source/public hashes, evaluates reviewer approval gates, builds
paragraph-level audit tables, and scores residual mosaic risk.

Visual proofing is parameterized separately from the release-audit text
path. Development renders enumerate four redaction treatments (blackout,
whiteout, grayout, blur) across four PDF backgrounds (white, gray,
black, blur) while preserving source-safe ledgers. The secure proof path
post-processes each proof PDF with the repository steganography layer:
hash manifests, provenance overlays, invisible text, barcodes, PDF
metadata, XMP metadata, embedded manifest attachments, and optional
Kmyth TPM \texttt{.ski} sidecars for the proof PDF and source-safe hash
manifest when runnable Kmyth tools and TPM access are available.

\newpage

\section{Results}\label{results}

The fixture release packet shows how invented high-side markers are
redacted before a public narrative is generated. The public export
contains redacted text only, paragraph audit rows, release authority
metadata, source-safe ledger hashes, review-gate status, and
residual-marker findings for review; it does not contain operational
source text beyond invented fixture residue used by tests.

The development proof matrix produces sixteen visual PDFs and, when
steganography is enabled, sixteen companion provenance PDFs with hash
sidecars. The matrix also records Kmyth runtime status and \texttt{.ski}
sidecar counts so TPM sealing is evidenced as an observed output rather
than a prose claim. This confirms that visual presentation choices
remain orthogonal to the release gate: the same source-safe decisions
drive every blackout, whiteout, grayout, blur, white-background,
gray-background, black-background, and blur-background output.



\clearpage

\bibliography{references}
\end{document}
